% ============================================================
% Section 39: 一维复式格子的晶格振动——双原子分子晶体的声子色散
% ============================================================
\newpage
\section{固体理论：一维复式格子的晶格振动——双原子分子晶体的声子色散}
\label{sec:diatomic-lattice-vibration}

\begin{modelbox}[题目：双原子分子晶体的简正振动模式]
讨论双原子分子晶体如固态 $\mathrm{H}_2$ 的简正振动模式。晶体可用一维复式格子模拟：格点上原子相同，质量为 $m$，相邻原子间距均为 $a/2$（晶格常数为 $a$），但最近邻原子间相互作用恢复力系数交错取值
\[
\beta_1 = 10c, \qquad \beta_2 = c.
\]
试求 $q=0$ 及布里渊区边界处的 $\omega(q)$，并定性画出色散关系。
\end{modelbox}

\begin{tipbox}[考点快速分析]
\begin{itemize}
    \item \textbf{一维复式格子}：基元含两个原子，出现声学支与光学支两支色散
    \item \textbf{运动方程与格波假设}：偶数、奇数原子位移耦合，得到 $2\times2$ 本征值问题
    \item \textbf{长波极限}：声学支 $\omega\propto q$（弹性波），光学支 $\omega(0)\neq 0$
    \item \textbf{布里渊区边界}：当 $\beta_1\neq\beta_2$ 时出现频率禁带
\end{itemize}
\end{tipbox}

\begin{derivationbox}[标准解题过程]
\textbf{Step 1: 模型与运动方程}

设第 $2n$ 个原子位移为 $u_{2n}$，第 $2n+1$ 个原子位移为 $u_{2n+1}$。运动方程为
\begin{align}
    m\ddot{u}_{2n} &= \beta_1(u_{2n+1}-u_{2n}) + \beta_2(u_{2n-1}-u_{2n}), \\
    m\ddot{u}_{2n+1} &= \beta_2(u_{2n+2}-u_{2n+1}) + \beta_1(u_{2n}-u_{2n+1}).
\end{align}

\textbf{Step 2: 格波假设与本征方程}

设格波解（注意奇数原子位置偏移 $a/2$）：
\[
u_{2n} = A e^{i(qna-\omega t)}, \qquad u_{2n+1} = B e^{i[q(n+\frac12)a-\omega t]}.
\]
代入运动方程，整理得关于 $A,B$ 的齐次线性方程组。非零解条件要求系数行列式为零，得到色散关系：
\begin{equation}
\boxed{m\omega^2 = (\beta_1+\beta_2) \pm \sqrt{\beta_1^2+\beta_2^2+2\beta_1\beta_2\cos(qa)}}
\end{equation}
或等价写成
\begin{equation}
\omega_{\pm}^2(q) = \frac{1}{m}\Bigl[(\beta_1+\beta_2) \pm \sqrt{(\beta_1+\beta_2)^2 - 4\beta_1\beta_2\sin^2\frac{qa}{2}}\Bigr].
\label{eq:diatomic-dispersion}
\end{equation}
其中 $\omega_-$ 为\textbf{声学支}，$\omega_+$ 为\textbf{光学支}。

\textbf{Step 3: 长波极限 $q\to 0$}

当 $q\to0$ 时，$\sin^2(qa/2)\approx (qa)^2/4$。对式 \eqref{eq:diatomic-dispersion} 中的根号作泰勒展开 $\sqrt{1-x}\approx 1-x/2$：

\textit{声学支}：
\begin{align}
    \omega_-^2(q) &\approx \frac{1}{m}\Bigl[(\beta_1+\beta_2) - (\beta_1+\beta_2)\Bigl(1-\frac{2\beta_1\beta_2}{(\beta_1+\beta_2)^2}\frac{(qa)^2}{4}\Bigr)\Bigr] \nonumber\\
    &= \frac{\beta_1\beta_2}{2m(\beta_1+\beta_2)}(qa)^2.
\end{align}
故
\begin{equation}
\boxed{\omega_-(q) \approx \sqrt{\frac{\beta_1\beta_2}{2m(\beta_1+\beta_2)}}\,a\,|q| \quad (q\to0)}
\end{equation}
即 $\omega_-\propto q$，对应弹性波；且 $\omega_-(0)=0$。

\textit{光学支}：
\begin{equation}
    \omega_+^2(q) \approx \frac{2(\beta_1+\beta_2)}{m} - \frac{\beta_1\beta_2}{2m(\beta_1+\beta_2)}(qa)^2.
\end{equation}
在 $q=0$ 处
\begin{equation}
\boxed{\omega_+(0) = \sqrt{\frac{2(\beta_1+\beta_2)}{m}} = \sqrt{\frac{22c}{m}}}
\end{equation}
（代入 $\beta_1=10c,\beta_2=c$）。

\textbf{Step 4: 布里渊区边界 $q=\pm\pi/a$}

在边界处 $\sin^2(qa/2)=1$，式 \eqref{eq:diatomic-dispersion} 简化为
\begin{equation}
    \omega_{\pm}^2\Bigl(\frac{\pi}{a}\Bigr) = \frac{1}{m}\Bigl[(\beta_1+\beta_2)\pm(\beta_1-\beta_2)\Bigr].
\end{equation}
因此
\begin{align}
    \omega_-^2\Bigl(\frac{\pi}{a}\Bigr) &= \frac{2\beta_2}{m} \quad\Longrightarrow\quad \boxed{\omega_-\Bigl(\frac{\pi}{a}\Bigr)=\sqrt{\frac{2c}{m}}}, \\[4pt]
    \omega_+^2\Bigl(\frac{\pi}{a}\Bigr) &= \frac{2\beta_1}{m} \quad\Longrightarrow\quad \boxed{\omega_+\Bigl(\frac{\pi}{a}\Bigr)=\sqrt{\frac{20c}{m}}}.
\end{align}
\end{derivationbox}

\begin{formulabox}[答案汇总]
\begin{center}
\begin{tabular}{ccc}
\toprule
波矢 $q$ & 声学支 $\omega_-$ & 光学支 $\omega_+$ \\
\midrule
$q=0$ & $0$ & $\displaystyle\sqrt{\frac{22c}{m}}$ \\[6pt]
$q=\pm\pi/a$ & $\displaystyle\sqrt{\frac{2c}{m}}$ & $\displaystyle\sqrt{\frac{20c}{m}}$ \\
\bottomrule
\end{tabular}
\end{center}
\end{formulabox}

\begin{insightbox}[色散关系的物理图像]
\begin{itemize}
    \item \textbf{声学支}：描述原胞质心的整体运动。$q=0$ 时所有原子同相振动，频率为零；向边界单调上升，斜率由等效弹性常数决定。
    \item \textbf{光学支}：描述原胞内两原子的相对振动。$q=0$ 时同胞内原子反相振动，频率最高；向边界略向下弯曲（因 $22c>20c$）。
    \item \textbf{频率禁带}：由于 $\beta_1\neq\beta_2$，布里渊区边界处两支频率不简并，存在禁带
    \[
    \Delta\omega = \sqrt{\frac{20c}{m}} - \sqrt{\frac{2c}{m}}.
    \]
    若 $\beta_1=\beta_2$，则禁带闭合，退化为单原子链的简并点。
    \item \textbf{与单缝/双缝的类比}（供记忆）：声学支像``基态''，光学支像``激发态''；禁带宽度直接反映两种键的强度差异。
\end{itemize}
\end{insightbox}

\begin{thinkbox}[考场速记：双原子链色散公式的快速重建]
核心结构：
\[
m\omega^2 = (\beta_1+\beta_2) \pm \sqrt{(\beta_1+\beta_2)^2 - 4\beta_1\beta_2\sin^2\frac{qa}{2}}.
\]
快速检验：
\begin{enumerate}
    \item $q=0$：根号内 $=(\beta_1+\beta_2)^2$，声学支 $m\omega^2=(\beta_1+\beta_2)-(\beta_1+\beta_2)=0$ \checkmark
    \item $q=\pi/a$：根号内 $=(\beta_1-\beta_2)^2$，两支分别为 $2\beta_2$ 与 $2\beta_1$ \checkmark
    \item 若 $\beta_1=\beta_2=\beta$：退化为单原子链 $m\omega^2 = 2\beta(1\pm\cos\frac{qa}{2})$，即 $4\beta\sin^2\frac{qa}{4}$ 与 $4\beta\cos^2\frac{qa}{4}$，边界处简并 \checkmark
\end{enumerate}
\end{thinkbox}
